% GPT-6 Astra Pro; /root/fixed_surface; 22 November 2026 UTC.
% Effort not exposed. Complete proof, submitted for independent review.
\section[HD-F1-S25: the intrinsic resolution and its actual deformation maps]{HD-F1-S25: canonical projectivization recovers the actual Abel resolution and identifies the normal-obstruction kernel}
\label{sec:f1-norm-fourfold-canonical-resolution}

Retain the actual component $a:T\to Y\subset P$ of S13--S24.
In particular $T$ is smooth, $Y$ is normal and Cohen--Macaulay,
$Ra_*\mathcal O_T=\mathcal O_Y$, and every nontrivial scheme
fibre is a complete pencil $E\simeq\mathbf P^1$ with
$N_{E/T}=\bigoplus_{i=1}^3\mathcal O_E(-b_i)$, where
$b_i\geq0$ and $\sum b_i=3$. The actual resolution is intrinsically
\[
 T\ \simeq\ \mathbf P_Y(\omega_Y)
       :=\operatorname{Proj}_Y\operatorname{Sym}_{\mathcal O_Y}(\omega_Y),
 \tag{F1S25.1}
\]
with the convention that projectivization represents invertible
quotients. The isomorphism is induced by the canonical evaluation
$a^*\omega_Y\twoheadrightarrow\omega_T$.

Write $D=\operatorname{Spec}\mathbf C[\epsilon]/(\epsilon^2)$.
Every actual embedded flat first-order deformation
$Y_D\subset P_D$ has the smooth proper lift
\[
 T_D=\mathbf P_{Y_D}(\omega_{Y_D/D})\longrightarrow Y_D\subset P_D.
 \tag{F1S25.2}
\]
Conversely, every deformation of the actual morphism $a$ with smooth
proper source $T_D$ and smooth ambient $P_D$ produces an embedded
flat deformation by its actual unit quotient
\[
 \mathcal O_{P_D}\twoheadrightarrow a_{D*}\mathcal O_{T_D}.
 \tag{F1S25.3}
\]
These constructions give the following exact description of the
kernel of the normal-cohomology map $\eta_Y$ of S24:
\[
 \ker\eta_Y
  =(a^*)^{-1}\operatorname{im}\left(
    H^1(T,T_T)\xrightarrow{H^1(da)}H^1(T,a^*T_P)\right).
 \tag{F1S25.4}
\]
Here $a^*:H^1(P,T_P)\to H^1(T,a^*T_P)$ is the actual pullback
isomorphism supplied by S21. This determines the precise tangent
map to be computed. It does not identify the deformations of $T$
with deformations of the original cyclic cover, and its image in
the sixteen-dimensional Weil tangent has not yet been evaluated.

\paragraph{1. The actual minors and their canonical projectivization.}
At a singular point use the complete local presentation constructed
in S24:
\[
 R=\mathbf C[[x_1,\ldots,x_6]],\qquad S=R/J,\qquad
 0\longrightarrow R^2\xrightarrow{\Phi}R^3
      \xrightarrow{f}R\longrightarrow S\longrightarrow0.
 \tag{F1S25.5}
\]
The entries of $f=(f_1,f_2,f_3)$ have initial forms the signed
minors $q$ of one of the three matrices $B_0$ in S24; the linear
part of $\Phi$ is $B_0^t$. Since $S$ is Cohen--Macaulay of
codimension two, finite duality gives
$\operatorname{Ext}_R^j(S,R)=0$ for $j\ne2$ and
\[
 \omega_S=\operatorname{coker}(\Phi^t:R^3\to R^2).
 \tag{F1S25.6}
\]
Let $c(\Phi)$ be the signed cofactor row of $\Phi$. The identity
$c(\Phi)\Phi=0$, together with exactness at the middle of the
dual of (F1S25.5), gives $c(\Phi)=u f$ for some $u\in R$.
Both rows have initial form $q$, so $u$ is a unit with constant
term one. Thus the actual ideal $J$ is generated by these actual
minors; this follows from the actual resolution rather than an
identification with the tangent cone.

Consequently $\mathbf P_S(\omega_S)$ is the closed subscheme of
$\operatorname{Spec}R\times\mathbf P^1$ defined by the three
equations
\[
                       \Phi\binom{s}{t}=0.
 \tag{F1S25.7}
\]
Indeed, on either projective chart the displayed linear relation
between the two columns makes all minors vanish, hence also the
original $f_i$. The other defining relations are exactly the
presentation (F1S25.6) of the symmetric algebra.
At a point over the closed point of $\operatorname{Spec}S$, the
derivative with respect to the six variables is the coefficient
matrix of $B_0^t(s,t)$. Its rank is three for every $[s:t]$.
For the three types this follows respectively from the independent
equations
\[
\begin{array}{ll}
 (111):& sa+td,\quad sb+te,\quad sc+tf;\\
 (012):& sx_0+tx_1,\quad sy_0+ty_1,\quad sy_1+ty_2;\\
 (003):& sy_0+ty_1,\quad sy_1+ty_2,\quad sy_2+ty_3.
\end{array}
 \tag{F1S25.8}
\]
In the second row the last two independent rows have rank two
and use variables disjoint from the first. In the third row one
of the two triangular three-by-three minors is $s^3$ or $t^3$.
The projectivization is therefore regular of dimension four at
every such point. Over the smooth locus of $Y$ its canonical
module is invertible, so projectivization is the identity. These
complete local checks prove that the actual global
$\mathbf P_Y(\omega_Y)$ is smooth.

\paragraph{2. Canonical evaluation identifies the given resolution.}
Proper duality for the given morphism $a$, normalized over
$\mathbf C$, supplies
\[
 Ra_*\omega_T[4]
  =R\mathcal Hom_Y(Ra_*\mathcal O_T,\omega_Y[4])
  =\omega_Y[4].
 \tag{F1S25.9}
\]
The resulting $a_*\omega_T=\omega_Y$ is the canonical trace
identification and is the identity on the common smooth open.
Adjunction on a pencil gives
\[
 \omega_T|_E=\omega_E\otimes\det N_{E/T}^{\vee}
             =\mathcal O_E(1).
 \tag{F1S25.10}
\]
S16 proves that the actual fibre ideal is
$\mathfrak m_y\mathcal O_T=I_E$. Its successive quotients are
$\operatorname{Sym}^mN_{E/T}^{\vee}$. Therefore
\[
 H^1\left(E,\omega_T|_E\otimes
                  \operatorname{Sym}^mN_{E/T}^{\vee}\right)=0
                       \qquad(m\geq0),
 \tag{F1S25.11}
\]
because each line summand has nonnegative degree, in fact degree
at least one. Applying the exact sequences of these actual
infinitesimal fibres and formal functions shows that the restriction
\[
 \widehat{(a_*\omega_T)_y}\longrightarrow
                            H^0(E,\mathcal O_E(1))
 \tag{F1S25.12}
\]
is surjective. Explicitly, every section lifts at every successive
order by (F1S25.11); choosing compatible lifts produces the
inverse-limit section. The target is annihilated by $\mathfrak m_y$,
so this is the ordinary fibre restriction map from $\omega_Y$.
Its source modulo $\mathfrak m_y$ has dimension two, by the
minimal presentation (F1S25.6). It is consequently an isomorphism
on these two-dimensional fibres.

Evaluation $a^*\omega_Y\to\omega_T$ is now onto along every
pencil: its restriction is the complete, globally generated
linear system of $\mathcal O_E(1)$, and Nakayama applies at
each point of $E$. It is an isomorphism away from the pencils.
This gives a morphism
$T\to\mathbf P_Y(\omega_Y)$ whose map on each exceptional
fibre is the complete $\mathcal O(1)$ isomorphism to $\mathbf P^1$.
It is the identity over the common smooth open. The morphism is
proper and has one-point fibres, hence is finite. It is surjective
on every fibre over $Y$, birational, and its target is smooth by
paragraph1. A finite birational morphism onto a normal integral
scheme is an isomorphism, proving (F1S25.1).

\paragraph{3. Every embedded first-order deformation has this smooth lift.}
Fix an actual embedded flat deformation $Y_D\subset P_D$.
We check its relative canonical module and (F1S25.2) in the
actual completed local rings. Lift the two independent linear
equations $z_1,z_2$ of S24 to the deformed ideal. This is possible
because the flat quotient makes reduction of that ideal onto
the central ideal surjective. The lifted equations still cut
out a smooth relative six-dimensional formal slice
$R_D=\mathbf C[\epsilon]/(\epsilon^2)[[x_1,\ldots,x_6]]$.
Write its quotient $S_D=R_D/J_D$. Both $S_D$ and $J_D$ are
flat over $D$.

Lift the three generators of $J$. Nakayama gives a surjection
$R_D^3\to J_D$. Its kernel is flat over $D$ and reduces to the
free rank-two kernel in (F1S25.5). Lifting a basis gives a
surjection $R_D^2$ onto that kernel. Its own kernel reduces to
zero, by flatness of the target, and is zero by Nakayama. Thus
we have the actual free resolution
\[
 0\longrightarrow R_D^2\xrightarrow{\Phi_D}R_D^3
       \xrightarrow{f_D}R_D\longrightarrow S_D\longrightarrow0.
 \tag{F1S25.13}
\]
To check its dual without presupposing a relative duality
base-change assertion, choose a regular sequence of length two
in the central ideal $J$ and lift it into $J_D$. Regularity,
and flatness over $D$ after quotienting by each member, lift
over the square-zero base: reduce a relation modulo $\epsilon$,
use the central nonzerodivisor, and then apply the same argument
to its $\epsilon$ coefficient. Hence $J_D$ contains such a
regular sequence and
$\operatorname{Ext}_{R_D}^j(S_D,R_D)=0$ for $j=0,1$.
For example the second vanishing follows by reducing along the
first member and using that the second is a nonzerodivisor in
the resulting ring. Dualizing (F1S25.13) therefore gives an exact
complex ending in
\[
 W_D=\operatorname{coker}(\Phi_D^t:R_D^3\to R_D^2)
        =\operatorname{Ext}_{R_D}^2(S_D,R_D).
 \tag{F1S25.14}
\]
The dual complex consists of $D$-flat modules and its reduction
is the exact central dual resolution. It computes
$\operatorname{Tor}_1^D(W_D,\mathbf C)=0$, so $W_D$ is
$D$-flat and reduces to $\omega_S$. Finite relative duality
for the actual quotient of the smooth relative slice identifies
$W_D$ with $\omega_{Y_D/D}$ in these coordinates. This also
proves, here, the required formation and base change of the
relative canonical sheaf. Changes of relative volume form tensor
the presentation by an invertible sheaf; its projectivization
and the intrinsic relative canonical sheaf glue in the usual
quotient convention.

As in paragraph1, the signed cofactor row of $\Phi_D$ is
$u_D f_D$: exactness of the dual complex proves the equality,
and reduction shows that $u_D$ is a unit. Thus the actual
projectivization over this slice is again defined by precisely
the three equations $\Phi_D(s,t)^t=0$. The Jacobian minors
checked in (F1S25.8) remain units at the corresponding points.
The relative Jacobian criterion proves smoothness of relative
dimension four over $D$. At smooth central points the canonical
module is invertible and the assertion follows from smoothness
of the flat deformation of that smooth locus. Every point of
the Artin-base deformation lies on the central topological
space, so these checks prove smoothness everywhere.
Projectivization is proper over the proper $Y_D$, and its
central fibre is the given $T$ by (F1S25.1). This proves the
actual lift (F1S25.2), functorially under isomorphisms of the
embedded deformations.

\paragraph{4. The reverse map is the actual flat unit quotient.}
Now start with a smooth proper deformation $T_D$ and a morphism
$a_D:T_D\to P_D$ reducing to the given $a$. Regard the central
objects as sheaves on the nilpotent thickenings. The exact
sequence on $T_D$ is
\[
 0\longrightarrow\mathcal O_T\xrightarrow{\epsilon}
       \mathcal O_{T_D}\longrightarrow\mathcal O_T\longrightarrow0.
 \tag{F1S25.15}
\]
The first arrow means the canonical identification of
$\epsilon\mathcal O_{T_D}$ with $\mathcal O_T$ supplied by
flatness. Applying the actual $Ra_{D*}$ and using S16 gives
\[
 R^qa_{D*}\mathcal O_{T_D}=0\ (q>0),\qquad
 0\longrightarrow\mathcal O_Y\longrightarrow F_D
        \longrightarrow\mathcal O_Y\longrightarrow0,
       \quad F_D=a_{D*}\mathcal O_{T_D}.
 \tag{F1S25.16}
\]
Multiplication by $\epsilon$ on $F_D$ is the displayed
surjection followed by the displayed inclusion. Thus
$\ker\epsilon=\epsilon F_D$, the flatness criterion over
$\mathbf C[\epsilon]/(\epsilon^2)$. In particular $F_D$ is
flat and $F_D/\epsilon F_D=\mathcal O_Y$; this is the actual
base-change map from (F1S25.15), not a presumed base-change
isomorphism.
Properness gives coherence of $F_D$. Its algebra unit
$\mathcal O_{P_D}\to F_D$ reduces to the original surjection
$\mathcal O_P\to\mathcal O_Y$. The coherent cokernel is
therefore zero by Nakayama. Its kernel defines the embedded
flat deformation $Y_D$, proving (F1S25.3). No assumption on
the origin of the deformation of $T$ has entered this argument.

\paragraph{5. The exact tangent-space comparison.}
For completeness the morphism-deformation map used here can
be read directly from actual gluing. Choose affine covers of
$T$ and $P$, refining the source so its members map into target
members. Trivialize first-order smooth deformations on these
affines. Their gluing differences are derivation cocycles
$\xi_{ij}$ for $T$ and $\kappa_{ij}$ for $P$. A local lift of
the map differs from the trivial lift by a section $b_i$ of
$a^*T_P$. The compatibility equation on overlaps is
\[
                 (\delta b)_{ij}=a^*\kappa_{ij}-da(\xi_{ij}).
 \tag{F1S25.17}
\]
This follows by composing the three actual first-order coordinate
maps and discarding $\epsilon^2$; it also constructs a global
morphism from a solution. Smooth affine schemes have no
first-order deformation obstructions on these local charts.
Using an acyclic affine refinement, (F1S25.17) says exactly
that $a^*\kappa$ belongs to the image of $H^1(da)$.

The normal map $\eta_Y$ is the embedded gluing obstruction in
the same convention: the coordinate derivation $v_{ij}$ changes
an equation $f\in I_Y$ by $v_{ij}(f)|_Y$. These differences
are the cocycle $\nu(v_{ij})$ of S24. A zero class is exactly
the existence of normal corrections to the local quotient
ideals making them glue. Equivalently, S23 compares the actual
quotient deformation with the perfect-object obstruction, and
the two injective maps in S24 identify that obstruction with
$\eta_Y$. Applying the two actual constructions in paragraphs3
and4 to (F1S25.17) proves (F1S25.4) in both directions.

Finally $T_P$ is trivial and S21 gives
$H^1(P,\mathcal O_P)\xrightarrow{a^*}H^1(T,\mathcal O_T)$
as an isomorphism. This proves the asserted actual pullback
isomorphism for $T_P$. Restricted to the actual polarized Weil
tangent $D_P$ of dimension sixteen, the image in question is
\[
 D_P\cap(a^*)^{-1}\operatorname{im}H^1(da).
 \tag{F1S25.18}
\]
It contains the twelve actual cover directions of S23. Its
dimension, and hence the rank of the remaining four-dimensional
map, remain to be computed.

\paragraph{6. The actual smooth-source normal sequence.}
The calculation can now be made on the given smooth fourfold
with its given embedding in $X=\operatorname{Sym}^8Z$:
\[
 0\longrightarrow T_T\longrightarrow T_X|_T
       \longrightarrow p^*\Omega_{\mathbf P^4}(1)
       \longrightarrow0.
 \tag{F1S25.19}
\]
This is the actual conormal computation of S13, dualized.
Let $u:X\to JZ$ be its Abel map, with the chosen constant
translation, and let $Q=2-\sigma-\sigma^2$ act on $JZ$.
The actual bundle map
\[
 T_X|_T\xrightarrow{du}\mathcal O_T\otimes\operatorname{Lie}(JZ)
      \xrightarrow{\frac13,dQ}
                   \mathcal O_T\otimes\operatorname{Lie}(P)
 \tag{F1S25.20}
\]
restricts to $da$ on $T_T$, because $dQ$ is multiplication
by three on $\operatorname{Lie}(P)$ and kills the invariant
summand. This is a map of complex vector bundles, so the
factor $1/3$ makes no assertion of an integral inverse isogeny.
Equations (F1S25.19)--(F1S25.20), followed by $H^1$, are an
explicit smooth-source description of the remaining map.

\paragraph{Primary sources and check scope.}
The proper-duality map in (F1S25.9) is the trace-compatible
isomorphism in Stacks, Lemmas 48.20.8--48.20.9 and
Example 48.20.11 (tags 0AUD, 0AUE, 0B6X), with
Example 48.3.9 (0GEW); their complete displayed proofs were
read. Both sources and targets are proper finite-type schemes
over $\mathbf C$, and their normalized dualizing complexes
are $\omega[4]$, as required. The finite local duality and
canonical-module input is the exact S24 input (0AX0,0AWS,
0B5A). Formal functions is applied to the actual fibre ideals
already proved in S16, with the vanishing checked explicitly
in (F1S25.11). The relative free resolution, its flatness and
base change, every Jacobian rank, the unit quotient and the
deformation gluing equation are proved above. No general
simultaneous-resolution theorem, formal cone identification,
or exhaustion by cover deformations is assumed.
